\subsection{Relations and Equivalence}

\begin{definition} \label{def:relations-and-equivalence} \hlt{Relations}\\
Let $X,Y$ be sets. A \hlt{relation} from $X$ to $Y$ is a subset $R\subseteq X\times Y$; write $xRy$ iff $(x,y)\in R$.
\end{definition}

\begin{definition} \label{def:equivalence-relation} \hlt{Equivalence Relations}\\
Let $X$ be a set. A relation $\sim$ on $X$ is an \hlt{equivalence relation} iff, for all $x,y,z\in X$,\\
$x\sim x$, $x\sim y\Rightarrow y\sim x$, and $(x\sim y\land y\sim z)\Rightarrow x\sim z$.
\end{definition}

\begin{definition} \label{def:equivalence-classes-invariants-canonical-forms} \hlt{Equivalence Classes, Invariants, and Canonical Forms}\\
Let $X,Y$ be sets.
\begin{enumerate}[label=(\roman*)]
\setlength{\itemsep}{0pt}
\item If $\sim$ is an equivalence relation on $X$, the equivalence class of $x\in X$ is $[x]_{\sim}\defeq\{y\in X:y\sim x\}$.
\item A map $f:X\to Y$ is an \hlt{invariant} of $\sim$ iff $x\sim y\Rightarrow f(x)=f(y)$ for all $x,y\in X$.
\item A map $f:X\to Y$ is a \hlt{complete invariant} of $\sim$ iff $x\sim y\Longleftrightarrow f(x)=f(y)$ for all $x,y\in X$.
\item A subset $\mathcal{C}\subseteq X$ is a \hlt{set of canonical forms} for $\sim$ iff $\abs{\mathcal{C}\cap[x]_{\sim}}=1$ for every $x\in X$.
\end{enumerate}
\end{definition}
